\begin{notechapter}
\chapter{Decomposition and Forecasting}\label{ch:decomposition-forecasting}
\chaptersubtitle{Trend, seasonality, and feeding a model its own predictions}

\begin{chapterintro}
Two practical questions come before any model of a time series. How do we separate a
slowly varying trend and a repeating seasonal pattern from the rest? And once a model
predicts one step ahead, how do we predict several steps ahead, when the
intermediate values are not yet observed? This chapter answers the first with the
classical decomposition, worked through on a series with multiplicative
seasonality, and the second by comparing recursive and direct forecasting,
including the two opposite meanings of ``open loop'' and ``closed loop''. You will
learn
\begin{itemize}[nosep]
  \item additive and multiplicative seasonality, and how a logarithm turns one into the
        other;
  \item the centred moving average, why its window must be centred, and what it leaves
        behind when the seasonality is multiplicative;
  \item how to estimate seasonal indices and remove them;
  \item recursive and direct multi-step forecasting, how forecast errors grow with the
        horizon, and the vocabulary of open and closed loops.
\end{itemize}
\prerequisites{\cref{ch:linear-time-series}.}
\end{chapterintro}

\section{Additive and multiplicative seasonality}\label{decomposition-forecasting:sec:models}

A classical decomposition writes the observation $Y_t$ as a combination of a \emph{trend} $T_t$
(slow changes of level), a \emph{seasonal} component $S_t$ of known period $s$, and an irregular
remainder $\varepsilon_t$:
\[
  \text{additive: } Y_t = T_t + S_t + \varepsilon_t,\qquad
  \text{multiplicative: } Y_t = T_t\times S_t\times\varepsilon_t .
\]
In the additive model the seasonal swings have a fixed size, and $S_t$ sums to $0$ over a
period. In the multiplicative model they are proportional to the level, and $S_t$ averages
to $1$ over a period: a value of $1.2$ means ``$20\%$ above trend''. For a positive series,
$\log Y_t = \log T_t + \log S_t + \log\varepsilon_t$ is additive, which is why multiplicative patterns are often
handled by modelling the logarithm.

\begin{example}[A series with multiplicative seasonality]\label{decomposition-forecasting:ex:series}
Monthly sales over three years, $t = 1,\dots,36$ ($t = 1$ is January of year 1), with a linear trend
and a sinusoidal seasonal factor, without noise:
\[
  T_t = 100 + 4t,\qquad S_t = 1 + 0.5\sin\frac{2\pi t}{12},\qquad Y_t = T_tS_t .
\]
$S_t$ has period $12$, averages to $1$, and swings by $\pm50\%$. The resulting values (one decimal) are
\begin{center}
\footnotesize\setlength{\tabcolsep}{3pt}
\begin{tabular}{@{}lcccccccccccc@{}}
\toprule
 & Jan & Feb & Mar & Apr & May & Jun & Jul & Aug & Sep & Oct & Nov & Dec \\
\midrule
year 1 & 130.0 & 154.8 & 168.0 & 166.2 & 150.0 & 124.0 & 96.0 & 74.8 & 68.0 & 79.4 & 108.0 & 148.0 \\
year 2 & 190.0 & 223.5 & 240.0 & 235.0 & 210.0 & 172.0 & 132.0 & 102.1 & 92.0 & 106.6 & 144.0 & 196.0 \\
year 3 & 250.0 & 292.3 & 312.0 & 303.8 & 270.0 & 220.0 & 168.0 & 129.3 & 116.0 & 133.8 & 180.0 & 244.0 \\
\bottomrule
\end{tabular}
\end{center}
The peak-to-trough range grows from $100$ in year 1 to $196$ in year 3: the signature of
multiplicative seasonality.
\end{example}

\section{Estimating the trend with a centred moving average}\label{decomposition-forecasting:sec:ma}

Averaging over exactly one period cancels a seasonal pattern that sums to zero over the
period. With a period of $12$, the plain average $\frac1{12}\sum_{k=-6}^5Y_{t+k}$ covers a full year, but its
window is not symmetric around $t$: it is centred at $t - \frac12$, between two months. The
\emph{centred} $2\times12$ moving average averages two consecutive such windows:
\[
  \mathrm{MA}_t = \frac1{24}\Bigl(Y_{t-6} + 2\sum_{k=-5}^5Y_{t+k} + Y_{t+6}\Bigr),
\]
a symmetric window of $13$ months whose end months get half weight. It is computed for
$t = 7,\dots,30$ (six values are lost at each end) by a convolution with the weights
$\frac1{24}(1, 2, 2,\dots,2, 1)$. Being symmetric, it reproduces a linear trend exactly, and it
annihilates any additive pattern of period $12$ that sums to zero.

\paragraph{What it leaves behind.} In a multiplicative series the seasonal effect is
$T_t(S_t - 1)$, whose size follows the trend, so the moving average does not cancel it
exactly. Writing $T_{t+k} = T_t + 4k$,
\[
  \mathrm{MA}_t = T_t\underbrace{\sum_kw_kS_{t+k}}_{=1} + 4\sum_kw_k\,k\,S_{t+k} = T_t + 4\sum_kw_k\,k\,(S_{t+k} - 1),
\]
using the symmetry of the weights. In the example the remainder oscillates between
$-3.7$ and $+3.7$ around the trend, about $1.5$ to $3\%$ of its level: small, but visible in
\cref{decomposition-forecasting:fig:series}.

\section{Seasonal indices}\label{decomposition-forecasting:sec:indices}

Dividing by the trend estimate isolates the seasonal and irregular parts:
$R_t = Y_t/\mathrm{MA}_t\approx S_t\varepsilon_t$ (for an additive model, subtract instead). Averaging $R_t$ over all years
for each calendar month, and rescaling the twelve averages to have mean $1$, gives the
\emph{seasonal indices}. In the example (two values per month are available):
\begin{center}
\footnotesize\setlength{\tabcolsep}{3pt}
\begin{tabular}{@{}lcccccccccccc@{}}
\toprule
 & Jan & Feb & Mar & Apr & May & Jun & Jul & Aug & Sep & Oct & Nov & Dec \\
\midrule
estimated & 1.227 & 1.418 & 1.500 & 1.448 & 1.272 & 1.020 & 0.767 & 0.574 & 0.500 & 0.561 & 0.736 & 0.978 \\
true $S_t$ & 1.250 & 1.433 & 1.500 & 1.433 & 1.250 & 1.000 & 0.750 & 0.567 & 0.500 & 0.567 & 0.750 & 1.000 \\
\bottomrule
\end{tabular}
\end{center}
The errors, at most $0.023$, come from the remainder of the moving average described
above. The \emph{seasonally adjusted} series is $Y_t/\hat S_t$; fitting a trend model to it and
multiplying back by $\hat S_t$ gives forecasts. Modern methods (STL, X-13ARIMA-SEATS)
refine this scheme with robust local regressions and seasonal patterns that evolve over
time.

\begin{figure}[ht]
\centering
\begin{tikzpicture}
\begin{groupplot}[group style={group size=1 by 2, vertical sep=1.4cm}, width=12cm, height=5cm,
  xmin=0, xmax=37, xtick={1,7,13,19,25,31}, xticklabels={Jan Y1, Jul Y1, Jan Y2, Jul Y2, Jan Y3, Jul Y3},
  tick label style={font=\scriptsize}, grid=major, grid style={gray!15}]
\nextgroupplot[ylabel={sales}, legend style={font=\scriptsize, draw=none, at={(0.02,0.98)}, anchor=north west}, legend cell align=left]
  \addplot[thick, black, mark=*, mark size=1pt] coordinates {(1,130.00) (2,154.77) (3,168.00) (4,166.23) (5,150.00) (6,124.00) (7,96.00) (8,74.84) (9,68.00) (10,79.38) (11,108.00) (12,148.00) (13,190.00) (14,223.55) (15,240.00) (16,235.01) (17,210.00) (18,172.00) (19,132.00) (20,102.06) (21,92.00) (22,106.59) (23,144.00) (24,196.00) (25,250.00) (26,292.33) (27,312.00) (28,303.80) (29,270.00) (30,220.00) (31,168.00) (32,129.27) (33,116.00) (34,133.81) (35,180.00) (36,244.00)};
  \addlegendentry{$Y_t$}
  \addplot[very thick, bookaccent, dashed] coordinates {(7,124.77) (8,130.13) (9,136.00) (10,141.87) (11,147.23) (12,151.73) (13,155.23) (14,157.87) (15,160.00) (16,162.13) (17,164.77) (18,168.27) (19,172.77) (20,178.13) (21,184.00) (22,189.87) (23,195.23) (24,199.73) (25,203.23) (26,205.87) (27,208.00) (28,210.13) (29,212.77) (30,216.27)};
  \addlegendentry{centred moving average}
  \addplot[thin, bookexample, domain=1:36] {100+4*x};
  \addlegendentry{true trend $100 + 4t$}
\nextgroupplot[ylabel={$Y_t/\mathrm{MA}_t$}, ymin=0.4, ymax=1.6]
  \addplot[thick, bookdef, mark=*, mark size=1pt] coordinates {(7,0.7694) (8,0.5751) (9,0.5000) (10,0.5595) (11,0.7335) (12,0.9754) (13,1.2240) (14,1.4161) (15,1.5000) (16,1.4495) (17,1.2745) (18,1.0222) (19,0.7640) (20,0.5729) (21,0.5000) (22,0.5614) (23,0.7376) (24,0.9813) (25,1.2301) (26,1.4200) (27,1.5000) (28,1.4457) (29,1.2690) (30,1.0173)};
  \addplot[thin, gray, domain=1:36, samples=200] {1+0.5*sin(deg(2*pi*x/12))};
\end{groupplot}
\end{tikzpicture}
\caption{Top: the series of \cref{decomposition-forecasting:ex:series}, its centred $2\times12$ moving
average, and the true trend; the moving average wobbles slightly around the trend
because the seasonality is multiplicative. Bottom: the ratio $Y_t/\mathrm{MA}_t$ (dots) against the
true seasonal factor $S_t$ (thin line).}
\label{decomposition-forecasting:fig:series}
\end{figure}

\section{Forecasting several steps ahead}\label{decomposition-forecasting:sec:multistep}

A model fitted to predict $X_{t+1}$ from the past, $\hat X_{t+1} = g(X_t, X_{t-1},\dots)$, can forecast $h$ steps ahead
in two ways.
\begin{itemize}[nosep]
  \item \emph{Recursive} (iterated) forecasting: predict $\hat X_{t+1}$, insert it in place of the unknown
        $X_{t+1}$, predict $\hat X_{t+2}$, and so on. One model serves all horizons.
  \item \emph{Direct} forecasting: fit a separate model $g_h$ for each horizon, $\hat X_{t+h} = g_h(X_t, X_{t-1},\dots)$,
        trained on pairs $(\text{past}, X_{t+h})$.
\end{itemize}
For a correctly specified linear model they coincide. For AR(1),
$\hat X_{t+h\mid t} = \alpha\hat X_{t+h-1\mid t} = \alpha^hX_t$, which is also $\E[X_{t+h}\mid X_t]$. The forecast error is
\[
  X_{t+h} - \alpha^hX_t = \sum_{j=0}^{h-1}\alpha^jZ_{t+h-j},\qquad
  \Var = \sigma_Z^2\,\frac{1 - \alpha^{2h}}{1 - \alpha^2}\ \xrightarrow[h\to\infty]{}\ \frac{\sigma_Z^2}{1 - \alpha^2} = \Var(X_t).
\]
With $\alpha = 0.8$ the error variance, in units of $\sigma_Z^2$, is $1$, $1.64$, $2.48$ and $2.75$ at horizons $1, 2, 5, 10$,
against $2.78$ for the unconditional variance: the forecast decays to the mean and its
error grows to the variance of the series. This growth is not a defect of the method; it
is the information lost about the future. What recursion adds is \emph{model} error: if the
one-step model is slightly wrong, its errors are fed back and compound, while a direct
model is fitted to the horizon it serves (at the cost of $h$ models, each with fewer useful
training pairs). For nonlinear models a further bias appears: iterating a point forecast
computes $g(\hat X_{t+1})$, whereas the conditional mean is $\E[g(X_{t+1})\mid\text{past}]$, and the two differ by
Jensen's inequality unless $g$ is linear.

\paragraph{Feeding back observations or predictions.} During forecasting, the inputs at
each step are either \emph{observed} values (when they arrive in time, as in a forecast
updated daily with new measurements) or the model's \emph{own predictions} (when they do
not, as in a ten-day forecast issued today). Using observations prevents error
accumulation but requires them; using predictions is the only option beyond one step.
The same distinction arises in training: \emph{teacher forcing} feeds the true previous values
to a sequence model during training, while at test time it must consume its own outputs.
The mismatch between the two regimes, called \emph{exposure bias}, is a source of
compounding errors in recurrent networks and language models.

\begin{pitfall}[``Open loop'' and ``closed loop'' have two opposite meanings]
In control engineering, a \emph{closed loop} feeds \emph{measured} outputs back into the
system, and an open loop runs without measurements; by analogy, closed-loop
forecasting uses observations as they arrive, and open-loop forecasting runs on the
model's predictions. In the neural-network time-series literature, in particular for
NARX networks (MATLAB's \texttt{narxnet}), the words are used the other way round: the
\emph{open-loop} (series--parallel) configuration feeds the \emph{true} past outputs to the
network, for training and one-step prediction, while the \emph{closed-loop} (parallel)
configuration feeds back the network's \emph{own predictions}, for multi-step
forecasting. Always check which convention a source uses; what matters is whether the
inputs at each step are observed values or predictions.
\end{pitfall}

\begin{chaptersummary}
\begin{itemize}[nosep]
  \item Additive seasonality has swings of fixed size, multiplicative seasonality swings
        proportional to the level; the logarithm turns the second into the first.
  \item The centred $2\times12$ moving average estimates the trend; it cancels additive seasonality
        exactly and leaves a small remainder for multiplicative seasonality with a trend.
  \item Seasonal indices are the per-month averages of $Y_t/\mathrm{MA}_t$ (or $Y_t - \mathrm{MA}_t$), normalized.
  \item Multi-step forecasts are recursive or direct. For AR(1), the $h$-step forecast is $\alpha^hX_t$, and its
        error variance $\sigma_Z^2(1 - \alpha^{2h})/(1 - \alpha^2)$ grows to the variance of the series.
  \item Whether a model is fed observations or its own predictions --- teacher forcing versus
        free running --- matters more than the name given to it, and the names
        ``open loop'' and ``closed loop'' are used in both senses.
\end{itemize}
\end{chaptersummary}

\section*{Exercises}
\addcontentsline{toc}{section}{Exercises}

\begin{exercise}\label{decomposition-forecasting:exr:1}
Show that the centred $2\times12$ moving average reproduces a linear trend exactly, and that it
annihilates any sequence $S_t$ of period $12$ with $\sum_{k=1}^{12}S_k = 0$.
\solutionlink{decomposition-forecasting:exr:1}
\end{exercise}

\begin{exercise}\label{decomposition-forecasting:exr:2}
In \cref{decomposition-forecasting:ex:series}, compute $\mathrm{MA}_7$ (July of year 1) and compare with $T_7 = 128$.
\solutionlink{decomposition-forecasting:exr:2}
\end{exercise}

\begin{exercise}\label{decomposition-forecasting:exr:3}
For the AR(2) model $X_t = X_{t-1} - 0.5X_{t-2} + Z_t$, write the recursive forecasts $\hat X_{t+1\mid t}$, $\hat X_{t+2\mid t}$ and
$\hat X_{t+3\mid t}$ in terms of $X_t$ and $X_{t-1}$, and the variance of the two-step forecast error.
\solutionlink{decomposition-forecasting:exr:3}
\end{exercise}
\end{notechapter}
